% sections/01_model.tex -- Section 1: the model, the priors, the sampler and the reported quantities.
% Owner: Writing Department, Task Lead 1.
% Sources (method v3.1.1 unchanged, plus the two reporting fields of the adopted evidence fix):
%   Docs/Writing/manuscript/spec/sections/sp1_notation.tex, sp2_model.tex, sp3_distributions.tex,
%   sp4_sampler.tex, sp4_algorithms.tex (spec of 2026-11-24);
%   Theory_work/Dept/prep/TH_PREP.md Sec. 1;
%   _local/engineering/reports/S1S2/perf/DECISION_MEMO.md (fix E) and designer_B/evidence_fix/README.md.
% No internal identifiers (tags, code names, field names) in the printed text.

\section{Model and method}\label{sec:model}

This section sets out the regression model, the priors and the posterior they give, the sequential Monte Carlo
sampler that SMC-BVS uses to explore the posterior, and the quantities a run reports. The model and priors are
standard. What is new is the sampler: it carries a population of models from a distribution that can be sampled
exactly to the posterior, and it builds that starting distribution from the correlation structure of the
predictors.

% ---------------------------------------------------------------------------------------------------------
\subsection{The linear model and the data}\label{sec:model-data}

The data $\Dn$ are a response $\yr\in\R^n$ and an $n\times\praw$ matrix $\Xr$ of predictors, measured on the same $n$
units. They follow the Gaussian linear model with an intercept,
\begin{equation}\label{eq:model}
  \yr=\alpha^*\one_n+\Xr\bstar+\varepsilon,\qquad \varepsilon\sim\Normal_n(0,\sigs^2I_n),
\end{equation}
where $\varepsilon$ is independent of $\Xr$. The number of predictors may be much larger than $n$, and the coefficient
vector $\bstar$ is sparse: the true model $\Sz:=\{j:\beta^*_j\neq0\}$ has a small size $\sz$. The aim is to recover
$\Sz$ and to quantify the uncertainty about it.

A predictor that is constant in the sample carries no information and is removed before anything else. From here on
$\Xr$ denotes the remaining $n\times p$ matrix, so $p$ always denotes the number of retained predictors, indexed
$1,\dots,p$ in their original order. All models below are sets of retained indices; the method reports them in the
original column numbers.

\paragraph{Standardisation.} The method centres and scales every column and the response. Let
$P^\perp_{\one}:=I_n-n^{-1}\one_n\one_n^\top$, let $D$ be the diagonal matrix of the sample standard deviations of
the columns of $\Xr$, and let $s_y$ be the sample standard deviation of $\yr$, all with divisor $n-1$. The method
works with
\begin{equation}\label{eq:standardise}
  X:=P^\perp_{\one}\Xr D^{-1},\qquad y:=P^\perp_{\one}\yr/s_y .
\end{equation}
Every column $x_j$ of $X$ then has mean zero and $\norm{x_j}^2=n-1$, and $y^\top y=n-1$. Centring removes the
intercept, which therefore has no prior. A coefficient $\theta_j$ and a variance $\sigma^2$ on the standardised
scale correspond to $\beta_j=s_y\theta_j/D_{jj}$ and $s_y^2\sigma^2$ in the original units. Section~\ref{sec:theory}
states its results for posterior draws $(S,\beta,\sigma^2)$ mapped back to the original units in this way, and for
the true coefficients $\bstar$ of~\eqref{eq:model}.

\paragraph{Models.} A model is a subset $S\subseteq\set{p}$, where $\set{p}:=\{1,\dots,p\}$. For a model $S$,
$X_S$ holds the columns of $X$ indexed by $S$, $G_S:=X_S^\top X_S$ is their Gram matrix and $z_S:=X_S^\top y$. The
method considers every model with at most $K$ predictors,
\begin{equation}\label{eq:space}
  \Om:=\{S\subseteq\set{p}:\abs{S}\le K\},\qquad K:=\min\big(p,\lfloor n/2\rfloor\big).
\end{equation}
The cap excludes only models with more than $n/2$ predictors, far larger than the sparse models of interest, and it
bounds every matrix the method factorises at $K\times K$. Given a model, the working likelihood on the standardised data is
\begin{equation}\label{eq:lik}
  y\mid S,\theta_S,\sigma^2\ \sim\ \Normal_n(X_S\theta_S,\ \sigma^2I_n),
\end{equation}
with coefficients $\theta_S\in\R^{\abs{S}}$ and variance $\sigma^2>0$.

% ---------------------------------------------------------------------------------------------------------
\subsection{Priors and posterior}\label{sec:model-prior}

\paragraph{Model prior.} Each predictor enters the model independently with probability $q$, and $q$ has a
$\BetaDist(1,p)$ prior. This is the beta-binomial prior $\BB(1,p)$, restricted to $\Om$. Integrating out $q$, a model
of size $k$ has prior mass proportional to
\begin{equation}\label{eq:modelprior}
  w(k):=\int_0^1q^k(1-q)^{p-k}\,p(1-q)^{p-1}\,dq=\frac{1}{2\binom{2p-1}{k}},\qquad
  \pi(S)=\frac{w(\abs{S})}{\sum_{k=0}^{K}\binom pk w(k)},\quad S\in\Om .
\end{equation}
The prior expected model size, before the cap, is $p/(p+1)$, close to one. Because $q$ is random, the prior odds of
adding one more predictor fall as $p$ grows, which is the multiplicity adjustment of \citet{ScottBerger2010}: moving
from size $k$ to size $k+1$ multiplies the prior mass of a model by $(k+1)/(2p-1-k)$.

\paragraph{Coefficient and variance priors.} Given $S$ and $\sigma^2$, the coefficients have Zellner's $g$-prior
\citep{Zellner1986} with $g=n$ and a unit ridge,
\begin{equation}\label{eq:slab}
  \theta_S\mid S,\sigma^2\ \sim\ \Normal_{\abs S}\big(0,\ n\sigma^2(G_S+I)^{-1}\big).
\end{equation}
With $g=n$ the prior carries about as much information as one observation \citep{KassWasserman1995}. The ridge
keeps the prior proper when the columns of $X_S$ are linearly dependent; since every diagonal entry of $G_S$ is $n-1$,
it barely changes the prior otherwise. The variance has the prior $\sigma^2\sim\IG(a_0,b_0)$ with $a_0=b_0=0.01$,
independent of $S$, on the standardised scale.

\paragraph{Posterior.} The priors are conjugate, so $\theta_S$ and $\sigma^2$ integrate out in closed form. Let
\begin{equation}\label{eq:HS}
  H_S:=(n+1)G_S+I,\qquad \SSq_S:=y^\top y-n\,z_S^\top H_S^{-1}z_S .
\end{equation}
Given $S$, the posterior of the parameters is
\begin{equation}\label{eq:condpost}
  \sigma^2\mid S,y\ \sim\ \IG\big(a_0+n/2,\ b_0+\SSq_S/2\big),\qquad
  \theta_S\mid S,\sigma^2,y\ \sim\ \Normal_{\abs S}\big(nH_S^{-1}z_S,\ n\sigma^2H_S^{-1}\big),
\end{equation}
and the marginal likelihood of $S$ is
\begin{equation}\label{eq:marglik}
  m(S)=C_n\Big\{\frac{\det(G_S+I)}{\det H_S}\Big\}^{1/2}\big(2b_0+\SSq_S\big)^{-(a_0+n/2)},\qquad
  C_n:=\frac{(2b_0)^{a_0}\,\Gamma(a_0+n/2)}{\pi^{n/2}\,\Gamma(a_0)},
\end{equation}
where $\pi=3.14\ldots$ in $C_n$; for the empty model the determinant ratio is one and $\SSq_\emptyset=y^\top y$.
The split-and-merge method of \citet{SongLiang2015} rests on the same conjugate structure, with a different slab
precision and an independent Bernoulli model prior under a size cap. We
call $F(S):=\log\pi(S)+\log m(S)$ the log score of $S$. The posterior over models and the evidence are
\begin{equation}\label{eq:posterior}
  \pis{S}=\frac{e^{F(S)}}{Z},\qquad Z:=\sum_{S\in\Om}e^{F(S)},\qquad S\in\Om .
\end{equation}
The posterior inclusion probability of predictor $j$ is $\PIP_j:=\sum_{S\ni j}\pis{S}$. The highest posterior model
$\HPM$ maximises $\pis{S}$, and the median probability model is $\MPM:=\{j:\PIP_j>1/2\}$ \citep{BarbieriBerger2004}.

The score shows the price a predictor pays to enter. If $x_j$ is orthogonal to the columns of $X_S$, adding $j$ to
a model of size $k$ changes $F$ by
\[
  -\log\frac{2p-1-k}{k+1}\;-\;\frac12\log n\;-\;\Big(a_0+\frac n2\Big)\log\frac{2b_0+\SSq_{S\cup\{j\}}}{2b_0+\SSq_S},
\]
a prior penalty of about $\log(2p)$, an Occam penalty of $\tfrac12\log n$, and a gain from the reduction in the
residual sum of squares. A predictor enters only if the last term outweighs the first two.

% ---------------------------------------------------------------------------------------------------------
\subsection{The sampler}\label{sec:model-sampler}

The model space has $\sum_{k\le K}\binom pk$ elements, far too many to enumerate when $p$ is large. How fast MCMC
samplers mix on this space depends on $p$ and on the correlation between predictors
\citep{YangWainwrightJordan2016,GriffinLatuszynskiSteel2021,ZhouYangVatsRobertsRosenthal2022}. SMC-BVS instead moves a population of models, called
particles, from a distribution that can be sampled exactly to the posterior, along a sequence of intermediate
distributions \citep{Chopin2002,DelMoralDoucetJasra2006}; \citet{SchaferChopin2013} use the same idea on binary
spaces. We describe the parts in the order in which a run uses them. Table~\ref{tab:defaults} lists the default
constants. The user sets only a random seed and a computing allocation.

\paragraph{Exact enumeration.} If $\abs{\Om}\le(5\Nc_0+4N_0)\,p$, where $N_0$ and $\Nc_0$ are the particle numbers of
Table~\ref{tab:defaults}, the method scores every model in $\Om$ and reports the exact posterior summaries. The
threshold equals the number of inclusion updates in a single round of the nine islands of a run (five in the
calibration, four in the final run; see below). Otherwise the method runs the sampler.

\paragraph{Blocks and the start law.} The predictors are partitioned into blocks $C_1,\dots,C_J$ of at most
$B:=\min(\lfloor\sqrt n\rfloor,\lfloor\sqrt p\rfloor,16)$ predictors each, by complete-linkage clustering with
distance $1-\abs{\hat r_{jl}}$, where $\hat r_{jl}$ is the sample correlation of predictors $j$ and $l$. Strongly
correlated predictors therefore tend to share a block. Within block $C_k$, every subset $U\subseteq C_k$ with
$\abs U\le K$ is scored, giving the table $h_k(U):=m(U)/m(\emptyset)$; a block has at most $2^{16}$ subsets. The
start law is
\begin{equation}\label{eq:startlaw}
  \pi_0(S)\ \propto\ \gamma_0(S):=\pi(S)\prod_{k=1}^{J}h_k(S\cap C_k),\qquad S\in\Om .
\end{equation}
It scores the predictors of each block jointly, by their exact marginal likelihood, and treats different blocks as
if they did not interact, apart from the prior on the total model size. Its
normalising constant $\bar Z_0:=\sum_{S\in\Om}\gamma_0(S)$ and independent exact draws from it are computed by a
recursion over blocks and model sizes, with no Monte Carlo error. The same recursion gives the mode of $\pi_0$, the
start mode $\tSz$.

\paragraph{The tempering path.} Let
\[
  \Delta(S):=\log\frac{m(S)}{m(\emptyset)}-\sum_{k=1}^{J}\log h_k(S\cap C_k),
\]
the part of the log marginal likelihood that the blocks do not explain; $\Delta(S)=0$ whenever $S$ lies within one
block. Since $e^{F(S)}=m(\emptyset)\,\gamma_0(S)\,e^{\Delta(S)}$, the distributions
\begin{equation}\label{eq:path}
  \pi_\xi(S)\ \propto\ \gamma_0(S)\,e^{\xi\Delta(S)},\qquad 0\le\xi\le1,
\end{equation}
run from the start law $\pi_0$ at $\xi=0$ to the posterior $\pis{\cdot}$ at $\xi=1$. The sampler moves along a grid
of temperatures $0=\grid0<\grid1<\dots<\grid{\Ts}=1$.

\paragraph{One step.} An island is a population of $N$ particles $S_1,\dots,S_N$, drawn independently from $\pi_0$.
The step from $\grid{t-1}$ to $\grid t$ has three parts:
\begin{enumerate}
  \item \emph{Reweight.} Particle $v$ receives the weight $\omega^{(t)}_v:=\exp\{(\grid t-\grid{t-1})\Delta(S_v)\}$.
  \item \emph{Resample.} $N$ particles are drawn with replacement in proportion to the weights (multinomial
    resampling, at every step).
  \item \emph{Move.} Each particle receives one round of Markov moves that leave $\pi_{\grid t}$ invariant.
\end{enumerate}
A round consists of $p+B+2$ moves applied in an order drawn uniformly at random for each particle:
\begin{itemize}
  \item an inclusion update of every predictor $j$: $j$ is added or deleted with its exact conditional probability
    under $\pi_\xi$ (a Gibbs step; an addition is not allowed at size $K$);
  \item one uniform swap: a member and a non-member of $S$, each chosen uniformly, are exchanged, with a
    Metropolis--Hastings acceptance step;
  \item $B$ neighbour swaps: a member $j$ of $S$ is exchanged for a non-member $l$ among the neighbours of $j$, with
    the Metropolis--Hastings ratio corrected for the asymmetric proposal;
  \item one independence move: a fresh draw $S'$ from $\pi_0$ is proposed, and it replaces $S$ with probability
    $\min\{1,e^{\xi\{\Delta(S')-\Delta(S)\}}\}$.
\end{itemize}
The neighbours $\Nb(j)$ of predictor $j$ are the at most $B$ predictors most correlated with $j$ among those with
$\abs{\hat r_{jl}}>\sqrt{2\log p/n}$, or the single most correlated predictor if none passes; the relation is then
made symmetric. A swap exchanges two predictors in one move, which inclusion updates can do only through an
intermediate model that may be much less probable; the independence move lets a particle jump to a model of a
different shape.

\paragraph{The schedule.} The grid is chosen in a separate calibration run and then frozen. A pilot island of $\Nc$
particles moves adaptively: from each temperature it takes the largest step for which the effective sample size of
the new weights, $\ESS=(\sum_v\omega_v)^2/\sum_v\omega_v^2$, is at least $0.90\,\Nc$, with a minimum step of
$1/\Tmax$. From the pilot's temperatures, a grid is built so that each step meets the same $0.90$ target, now
computed exactly on a candidate set: the empty model and the $20p$ best models found so far by $\gamma_0$ and by
$F$. Four diagnostic islands are run on this grid, the candidate set grows, and the grid is rebuilt, up to $\ell_n$
times or until it stops changing. If any of seven calibration warnings fires (an island's best score more than one
unit below the best found, a standard deviation above one among the islands' log evidence estimates, a local search
that improves the best model, a local search that stops at its move cap, a forced minimum step, a failed grid, an
unconverged grid), the calibration is \emph{escalated}: it is repeated once, with four times as many particles and a
four times higher step ceiling, and the final run then uses four times as many particles too. A grid that fails to
be built even then is replaced by the last grid built successfully, or by the pilot's temperatures. The frozen grid does not depend on the
random numbers of the final run.

\paragraph{The final run.} Four islands, each of $N$ particles and with its own random stream, run on the frozen
grid. At $\xi=1$ each island applies $\Lbi$ further rounds as burn-in, then four measured rounds, from which the
inclusion probabilities are estimated. Island $r$ estimates the evidence by
\begin{equation}\label{eq:zhat}
  \Zhat_r:=m(\emptyset)\,\bar Z_0\prod_{t=1}^{\Ts}\Big(\frac1N\sum_{v=1}^N\omega^{(t)}_v\Big),
\end{equation}
with $\omega^{(t)}_v$ computed on the particles before resampling at step $t$. Given the frozen calibration, and for a run that
is not stopped by the time or memory limit, $\Zhat_r$ is an unbiased estimator of $Z$, a standard property of this
construction \citep{DelMoral2004}; $\log\Zhat_r$ is not unbiased for $\log Z$. The pooled estimate is
$\Zhat:=\frac14\sum_{r=1}^4\Zhat_r$; the islands are weighted equally.

\paragraph{Local searches.} After the islands, two deterministic searches over the neighbourhood of a model (all
models one deletion, one addition or one swap away) climb the log score. At each move the search scans the whole
neighbourhood and moves to its best member, if that member strictly improves $F$. The first search starts from the
best model scored so far, with at most $\ell_n$ moves; the second starts from the start mode $\tSz$, with at
most $\ell_n+\abs{\tSz}$ moves. A search is \emph{certified} when it stops because no neighbour has a strictly higher
score.

\paragraph{Resources.} Before it starts, the method checks its memory and time needs against the allocation (by
default 8\,GiB and 24\,h). A run that would exceed the allocation, runs out of time, or meets a numerical failure in
scoring stops and reports no estimates.

\begin{table}[t]
\centering\small
\caption{Default constants of SMC-BVS. Here $i\in\{0,1\}$ is the calibration tier: $i=1$ only after a warning in
the calibration at $i=0$.}\label{tab:defaults}
\begin{tabular}{@{}lll@{}}
\toprule
Quantity & Symbol & Default \\
\midrule
model-size cap & $K$ & $\min(p,\lfloor n/2\rfloor)$ \\
block size cap & $B$ & $\min(\lfloor\sqrt n\rfloor,\lfloor\sqrt p\rfloor,16)$ \\
particles per island, final run & $N=N_i$ & $4^iN_0$, with $N_0:=\max(2500,n,p)$ \\
particles per island, calibration & $\Nc=\Nc_i$ & $4^i\lceil512N_0/2500\rceil$ \\
step ceiling & $\Tmax=\Tmax_i$ & $4^i\max(1000,\ell_n^2)$ \\
move cap of the local search & $\ell_n$ & $\lceil\log_2(n+p+1)\rceil$ \\
burn-in rounds at $\xi=1$ & $\Lbi$ & $\max(4,\lceil\log_2(n+1)\rceil)$ \\
islands; measured rounds & & $4$; $4$ \\
effective sample size target & & $0.90$ \\
returned list of top models & $\Tlist$ & the best $40p+1$ models scored \\
\bottomrule
\end{tabular}
\end{table}

% ---------------------------------------------------------------------------------------------------------
\subsection{Variable selection and the reported quantities}\label{sec:model-output}

A run ends in one of three ways: \emph{exact} (the model space was enumerated), \emph{completed} (the sampler ran to
the end) or \emph{stopped}. A stopped run reports no estimates. A completed run reports the following.

\paragraph{Inclusion probabilities and the reported model.} Within island $r$, $\PIPhat_{r,j}$ is the average, over
the $N$ particles and the four measured rounds, of the conditional inclusion probability computed in the inclusion
update of predictor $j$. Averaging these probabilities rather than the inclusion indicators reduces the Monte Carlo
noise. The method reports
\begin{equation}\label{eq:pipmpm}
  \PIPhat_j:=\frac14\sum_{r=1}^4\PIPhat_{r,j},\qquad \MPMhat:=\{j:\PIPhat_j>1/2\}.
\end{equation}
The estimated median probability model $\MPMhat$ is the model SMC-BVS reports. When we compare methods, $\Shat$
denotes the model a method reports, so $\Shat=\MPMhat$ for SMC-BVS.

\paragraph{The best model found.} The method also stores $\HPMhat$, the model with the highest log score among all
models scored during the run: in the block tables, the calibration, the islands and both local searches. Ties go to
the smaller model, then to the lexicographically smaller one. $\HPMhat$ is \emph{certified} when it is the end point
of a certified local search; it is then a local maximum of $F$ over its neighbourhood. $\HPMhat$ estimates the
$\HPM$; it is not the reported model.

\paragraph{The list of top models.} The method returns the list $\Tlist$ of the best models it scored, at most
$40p+1$ of them, with their log scores. Renormalising the posterior over the list gives
\begin{equation}\label{eq:list}
  \Zlist:=\sum_{S\in\Tlist}e^{F(S)},\qquad
  \piplist(j):=\frac{1}{\Zlist}\sum_{S\in\Tlist,\,S\ni j}e^{F(S)},\qquad
  \mpmlist:=\{j:\piplist(j)>1/2\}.
\end{equation}
These are reported beside $\PIPhat$ and $\MPMhat$. Since $\Tlist\subseteq\Om$, the list evidence satisfies
$\Zlist\le Z$, with no Monte Carlo error.

\paragraph{The evidence.} The Monte Carlo estimate $\Zhat$ and the bound $\Zlist$ are combined. The reported log
evidence and its lower limit are
\begin{equation}\label{eq:evidence}
  \log Z_{\mathrm{rep}}:=\max\big(\log\Zhat,\ \log\Zlist\big),\qquad
  \log Z_{\mathrm{low}}:=\max\big(\log(0.05\,\Zhat),\ \log\Zlist\big).
\end{equation}
Taking the maximum never moves the point estimate away from $Z$, because $\Zlist\le Z$; it can bias the estimate
upward on the $Z$ scale, so the raw $\Zhat$ remains the unbiased estimate. Because $\Zhat$ is unbiased (under the conditions stated in Section~\ref{sec:model-sampler}), Markov's
inequality gives $\P(0.05\,\Zhat>Z)\le0.05$, and since $\Zlist\le Z$ the lower limit
$Z_{\mathrm{low}}$ exceeds $Z$ on exactly the same event. It is therefore a 95\% one-sided lower limit for $Z$. The
raw $\log\Zhat$, the four island values $\log\Zhat_r$ and the raw limit are reported too, together with the gap
$\log\Zlist-\log\Zhat$ and a flag when it exceeds one: $\Zhat$ then underestimates $Z$ by more than a factor $e$.

\paragraph{Exact runs.} An exact run reports the exact $\PIP_j$, $\MPM$ and $\log Z$, and $\MPM$ is the reported
model; the best model found is the $\HPM$ and is certified. The list fields are computed from the $40p+1$ best models, and the lower limit is
$\max\{\log(0.05\,Z),\log\Zlist\}$.

\paragraph{Diagnostics.} After the final run, four checks are recomputed as flags; they never change an output:
(i) some island's best score lies more than one unit below $F(\HPMhat)$; (ii) the island log evidence estimates
$\log\Zhat_r$ have standard deviation above one; (iii) the local search from the best model scored so far raised the
best score (measured before the search from the start mode); (iv) the effective sample size fell below $N/2$ at
some step. We say that a run \emph{raised a run-time warning} when at least one of these four flags is set. A run
also records, without counting them as run-time warnings, whether the calibration was escalated and which
calibration warnings fired, whether either local search stopped at its move cap, the models of $\Tlist$ within
$10^{-9}$ of the best score, and the pairs of predictors with $\abs{\hat r_{jl}}\ge1-10^{-10}$.
